% Extended Abstract --- 1-page grading-critical artifact (CS285 final
% project rubric, p.2). Five paragraphs: setup+headline,
% basin-selection mechanism revision (lifts the load-bearing-function
% sentence from Sec.~\ref{sec:discussion} verbatim), failure-mode
% landscape, methodological contributions, implications + scope.
% All numbers cross-checked against results/paper_master_table.csv.
% Compressed in v2 to fit on 1 page in A4+1in geometry.

\section*{Extended Abstract}

We study reinforcement learning for portfolio allocation across eight
broad-market ETFs at daily frequency: the agent emits weights on the
simplex and the env rewards the post-cost log return $r_t = \log(1 +
w_t^\top R_{t+1} - \lambda \|w_t - w_{t-1}\|_1)$. Our milestone
established that the standard offline-RL stack collapses smoothly
under realistic transaction-cost friction on this universe:
offline-trained policies fail to reach the equal-weight baseline at
$\lambda{=}0.001$ and actively destroy capital at $\lambda{=}0.005$.
This paper asks whether offline-to-online (O2O) fine-tuning with
adaptive conservatism scheduling repairs the collapse. The headline
empirical result, on a 1{,}061-day causal test window
(2022-Q1 to 2026-Q1, the post-2022 stock-bond-correlation regime), is
yes: an adaptive-O2O agent (Geodesic-CQL pretrain on 2008--2020
transitions, SAC-Dirichlet fine-tune with $\beta_t = \alpha_{\mathrm{CQL}}
\cdot \sigma(\eta \cdot \mathrm{KL}(h_{\text{off}} \| h_{\text{on}}))$)
achieves test Sharpe $1.762 \pm 0.154$ across three seeds --- the only
RL method in our study to beat both equal-weight ($0.953$) and
risk-parity ($1.161$) on this window.

The mechanism behind that headline is not what we pre-registered. Our
pre-registered hypothesis was that the modulated $\beta_t$ would let
the SAC fine-tune trade actively in response to distribution shift in
the test window. The Phase~2C results disconfirm that mechanism. Both naive ($\beta_t \equiv
\alpha_{\mathrm{CQL}}$) and adaptive O2O converge to functionally
static allocations (per-seed turnover $\le 2.2\!\times\!10^{-4}$
across all six runs, three orders of magnitude below GRPO online's
$0.24$); the differential between them is not how they trade but which
static allocation the SAC fine-tune lands on. Naive O2O's pinned
conservatism leaves the basin random across seeds (Sharpe range
$[-0.236,\,+2.179]$, std $1.21$); adaptive O2O's modulated schedule
consistently steers the same SAC dynamics from the same Geodesic-CQL
pretrain to a uniformly strong basin (Sharpe range $[+1.629,\,+1.931]$,
std $0.15$, an eight-fold variance reduction). The conservatism
schedule's load-bearing function in our results is convergence-basin
steering during fine-tune, not online responsiveness during deployment.

The four-way O2O comparison sits inside a broader failure-mode
landscape, with each mode mechanistically traceable. (i) On a leaky
offline pipeline (in-observation reward feature, since fixed;
Appendix~\ref{app:leak}), $Q$-maximizing offline RL exploits the leak:
TD3+BC test Sharpe $6.785 \pm 0.628$, BCQ $3.905 \pm 0.232$, vs
behavior-anchored algorithms at $\sim\!0.94$. (ii) SAC-Dirichlet's auto-temperature blows up ($\tau \sim 10^9$);
the online-only policy collapses to a near-equal-weight static
allocation identical across seeds. (iii) Naive O2O exhibits the
seed-dependent-static pathology unpacked above. (iv) Online GRPO
actively trades (turnover $0.24$) at net Sharpe cost ($0.297 \pm
0.017$); even the best group-size ablation cell ($G{=}4$, Sharpe
$0.500 \pm 0.108$) lands below equal-weight.

Three methodological contributions. \emph{Continuous GRPO}: a critic-free adaptation of
Group Relative Policy Optimization to a continuous Dirichlet actor on
the simplex. The derivation exploits next-state independence from the
agent's action (\emph{exogeneity}; unit-tested at the env boundary),
under which the bootstrapped continuation cancels under within-group
centering. The formulation is sound; whether it beats classical
baselines is a separate empirical question, settled here in the
negative.
\emph{Algorithm-class differential as a leak diagnostic}: a $> 2\times$
Sharpe gap between unconstrained $Q$-max and conservative-$Q$/behavior-anchored
methods on the same dataset signals label leakage (CQL\_vanilla as the
``no signal''/``leakage'' control row).
\emph{Scope condition}: the behavior-anchored basin is
feature-space-specific --- on a 216-d feature space, IQL exits the
basin and ends in active trading (Sharpe $0.40$, turnover $0.21$, vs.
the 56-d basin's $\sim\!0.94$).

The basin-selection finding suggests an underweighted axis for
conservatism scheduling research --- seed-stability of the fine-tune's
converged fixed point under weakly identified dynamics, not just
online responsiveness to drift --- evaluated here under one universe,
one $\lambda$, one behavior mixture, one test window.
